\documentclass[12pt,a4paper]{article}

% -------------------------------------------------------
% Paquetes
% -------------------------------------------------------
\usepackage[utf8]{inputenc}
\usepackage[T1]{fontenc}
\usepackage[spanish]{babel}
\usepackage{geometry}
\geometry{left=2.5cm, right=2.5cm, top=3cm, bottom=3cm}
\usepackage{graphicx}
\usepackage{hyperref}
\usepackage{listings}
\usepackage{xcolor}
\usepackage{booktabs}
\usepackage{float}
\usepackage{amsmath}
\usepackage{enumitem}
\usepackage{titlesec}
\usepackage{parskip}
\usepackage{fancyhdr}
\usepackage{caption}

% -------------------------------------------------------
% Colores para código
% -------------------------------------------------------
\definecolor{codebg}{RGB}{245,245,245}
\definecolor{codegreen}{RGB}{0,128,0}
\definecolor{codegray}{RGB}{100,100,100}
\definecolor{codeblue}{RGB}{0,0,180}
\definecolor{codered}{RGB}{180,0,0}

\lstdefinestyle{python}{
  backgroundcolor=\color{codebg},
  basicstyle=\ttfamily\footnotesize,
  keywordstyle=\color{codeblue}\bfseries,
  commentstyle=\color{codegreen},
  stringstyle=\color{codered},
  numberstyle=\tiny\color{codegray},
  breaklines=true,
  captionpos=b,
  numbers=left,
  frame=single,
  language=Python
}

\lstdefinestyle{yaml}{
  backgroundcolor=\color{codebg},
  basicstyle=\ttfamily\footnotesize,
  keywordstyle=\color{codeblue},
  commentstyle=\color{codegreen},
  stringstyle=\color{codered},
  breaklines=true,
  captionpos=b,
  frame=single
}

% -------------------------------------------------------
% Encabezado y pie de página
% -------------------------------------------------------
\pagestyle{fancy}
\fancyhf{}
\rhead{Secure DevOps Pipeline}
\lhead{ESPE -- Desarrollo de Software Seguro}
\cfoot{\thepage}

% -------------------------------------------------------
% Datos del documento
% -------------------------------------------------------
\title{
  \vspace{-1cm}
  \includegraphics[width=3cm]{espe_logo.png}\\[0.5cm]
  \textbf{Informe Técnico}\\[0.3cm]
  \large Desarrollo e Implementación de un Pipeline CI/CD Seguro\\
  con Detección Automática de Vulnerabilidades mediante\\
  un Modelo de Minería de Datos\\[0.5cm]
  \normalsize Proyecto Integrador Parcial II
}
\author{
  Gabriel Nicolás Vivanco Raza\\
  \small Ingeniería en Software\\
  \small Desarrollo de Software Seguro\\
  \small Profesor: Geovanny Cudco
}
\date{Junio 2026}

% -------------------------------------------------------
\begin{document}
% -------------------------------------------------------

\maketitle
\thispagestyle{empty}

\begin{abstract}
Se presenta el diseño e implementación de un pipeline de Integración y Entrega Continua (CI/CD) completamente automatizado, orientado a la seguridad del software. El sistema integra un clasificador binario XGBoost —entrenado sobre un dataset público de 35\,000 fragmentos de código Java— capaz de determinar si el código fuente enviado a revisión es \textit{SEGURO} o \textit{VULNERABLE}. El pipeline ejecuta tres etapas secuenciales y obligatorias: (1)~análisis de seguridad mediante ML, (2)~ejecución de pruebas unitarias y de integración, y (3)~despliegue automático en producción. Se implementan notificaciones en tiempo real a través de un bot de Telegram y se generan \textit{issues} automáticas en GitHub ante cualquier fallo. La aplicación se encuentra desplegada en Render y accesible de forma pública.
\end{abstract}

\newpage
\tableofcontents
\newpage

% -------------------------------------------------------
\section{Introducción}
% -------------------------------------------------------

El paradigma \textit{Shift-Left Security} propone integrar controles de seguridad en las fases más tempranas del ciclo de vida del software. En la práctica, esto implica que la revisión de vulnerabilidades no debe limitarse a auditorías manuales realizadas después del desarrollo, sino que debe formar parte integral del pipeline automatizado de CI/CD.

El presente proyecto implementa esta filosofía mediante la construcción de un sistema que analiza el código fuente Java en el momento en que el desarrollador crea un \textit{Pull Request}, antes de que dicho código llegue siquiera a la rama de pruebas. El análisis se delega a un modelo de aprendizaje automático clásico —\textbf{XGBoost}— entrenado con ejemplos reales de código vulnerable y seguro, garantizando la objetividad y reproducibilidad del proceso.

\textbf{Objetivo principal:} Diseñar, implementar y demostrar un pipeline CI/CD completamente automatizado y seguro que integre un modelo de inteligencia artificial basado en técnicas de minería de datos, capaz de clasificar código fuente como SEGURO o VULNERABLE, permitiendo que únicamente el código considerado seguro llegue a producción.

\textbf{Nota importante:} El proyecto utiliza exclusivamente modelos de clasificación de minería de datos clásicos (XGBoost, scikit-learn). No se emplean modelos de lenguaje a gran escala (LLMs) como GPT, Claude, Llama ni ningún derivado.

% -------------------------------------------------------
\section{Arquitectura del Sistema}
% -------------------------------------------------------

\subsection{Visión general}

El sistema está compuesto por tres capas funcionales:

\begin{enumerate}
  \item \textbf{Capa de análisis de seguridad:} módulo Python que extrae 14 características léxicas y estructurales del código Java y las procesa con el modelo XGBoost.
  \item \textbf{Capa de orquestación CI/CD:} flujo de trabajo en GitHub Actions con tres \textit{jobs} secuenciales.
  \item \textbf{Capa de aplicación:} API REST Flask que expone el clasificador, y una aplicación Spring Boot que sirve como código objetivo.
\end{enumerate}

\subsection{Ramas del repositorio}

\begin{table}[H]
\centering
\begin{tabular}{@{}lll@{}}
\toprule
\textbf{Rama} & \textbf{Propósito} & \textbf{Protección} \\
\midrule
\texttt{dev}  & Desarrollo activo (push del desarrollador) & Sin restricciones \\
\texttt{test} & Staging / pruebas                          & Requiere checks de CI \\
\texttt{main} & Producción                                 & Requiere checks de CI \\
\bottomrule
\end{tabular}
\caption{Ramas del repositorio y sus roles}
\end{table}

\subsection{Flujo de trabajo}

El pipeline se activa automáticamente al crear un Pull Request de \texttt{dev} hacia \texttt{test}. El flujo completo es el siguiente:

\begin{enumerate}
  \item El desarrollador realiza un push a \texttt{dev} y abre un PR hacia \texttt{test}.
  \item GitHub Actions ejecuta el \textit{job} \texttt{security-check}: descarga el diff del PR, extrae las features del código Java modificado, predice con XGBoost y decide si continuar o bloquear.
  \item Si el código es SEGURO, se ejecuta el \textit{job} \texttt{tests}: pruebas Python (pytest) y Java (JUnit 5 + Maven).
  \item Al producirse el push a \texttt{test} (resultado del merge automático), el \textit{job} \texttt{merge-and-deploy} fusiona \texttt{test} $\to$ \texttt{main} y activa el webhook de despliegue en Render.
\end{enumerate}

% -------------------------------------------------------
\section{Modelo de Minería de Datos}
% -------------------------------------------------------

\subsection{Algoritmo seleccionado: XGBoost}

Se seleccionó XGBoost (\textit{Extreme Gradient Boosting}) como clasificador principal por las siguientes razones:

\begin{itemize}
  \item Excelente rendimiento en datos tabulares con features de conteo.
  \item Robustez ante \textit{outliers} sin preprocesamiento exhaustivo.
  \item Capacidad de capturar interacciones no lineales entre características.
  \item Soporte nativo para clasificación binaria con salida de probabilidades.
\end{itemize}

\subsection{Configuración del modelo}

\begin{lstlisting}[style=python, caption={Configuración del clasificador XGBoost}]
model = XGBClassifier(
    n_estimators=200,
    max_depth=8,
    learning_rate=0.1,
    subsample=0.8,
    colsample_bytree=0.8,
    random_state=42,
    use_label_encoder=False,
    eval_metric='logloss',
    n_jobs=-1
)
\end{lstlisting}

\subsection{Pipeline de entrenamiento}

\begin{enumerate}
  \item Carga del dataset \texttt{data/processed/security\_dataset.csv} (columnas \texttt{code}, \texttt{label}).
  \item Extracción de las 14 características por cada registro mediante \texttt{JavaCodeAnalyzer}.
  \item Normalización con \texttt{StandardScaler} (ajustado únicamente sobre el conjunto de entrenamiento).
  \item División estratificada 80\% entrenamiento / 20\% prueba (\texttt{random\_state=42}).
  \item Entrenamiento del clasificador XGBoost con conjunto de validación interno.
  \item Evaluación: accuracy, precision, recall, F1-score y validación cruzada 5-fold.
  \item Serialización del modelo (\texttt{classifier\_model.joblib}) y del scaler (\texttt{scaler.joblib}).
\end{enumerate}

\subsection{Dataset}

\begin{table}[H]
\centering
\begin{tabular}{@{}ll@{}}
\toprule
\textbf{Atributo}        & \textbf{Valor} \\
\midrule
Fuente                   & Dataset público de vulnerabilidades (Vulnerability Fix Dataset) \\
Registros originales     & 35,000 \\
Columnas originales      & \texttt{vulnerability\_type}, \texttt{vulnerable\_code}, \texttt{fixed\_code} \\
Tipos de vulnerabilidad  & SQL Injection, Command Injection, XSS, \\
                         & Insecure Deserialization, Path Traversal \\
Registros duplicados eliminados  & 2,517 (7.19\%) \\
Registros texto-explicativo elim. & 11,395 (fixed\_code en inglés) \\
Dataset final            & 42,176 filas — perfectamente balanceado (50/50) \\
Formato final            & \texttt{code}, \texttt{label} (0 = SEGURO, 1 = VULNERABLE) \\
\bottomrule
\end{tabular}
\caption{Características del dataset de entrenamiento}
\end{table}

\subsection{Extracción de características}

La clase \texttt{JavaCodeAnalyzer} (archivo \texttt{security\_checker/feature\_extractor.py}) implementa la extracción de 14 características numéricas a partir del texto del código fuente:

\begin{table}[H]
\centering
\begin{tabular}{@{}lll@{}}
\toprule
\textbf{Feature}               & \textbf{Categoría}    & \textbf{Descripción} \\
\midrule
\texttt{code\_length}          & Estructural           & Total de caracteres (log) \\
\texttt{line\_count}           & Estructural           & Total de líneas \\
\texttt{token\_count}          & Léxica                & Total de tokens léxicos (log) \\
\texttt{keyword\_count}        & Léxica                & Palabras clave Java \\
\texttt{comment\_ratio}        & Estilo                & Proporción de comentarios \\
\texttt{bracket\_depth}        & Complejidad           & Profundidad máxima de llaves \\
\texttt{brace\_count}          & Complejidad           & Número de llaves \texttt{\{ \}} (log) \\
\texttt{parenthesis\_count}    & Complejidad           & Número de paréntesis (log) \\
\texttt{dangerous\_function\_count} & Seguridad        & Llamadas a APIs peligrosas \\
\texttt{sql\_pattern\_count}   & Seguridad             & Patrones de SQL injection \\
\texttt{system\_call\_count}   & Seguridad             & Llamadas a \texttt{System.} \\
\texttt{reflection\_count}     & Seguridad             & Uso de reflexión Java \\
\texttt{sanitization\_count}   & Protección            & Mecanismos de sanitización \\
\texttt{try\_catch\_count}     & Manejo de errores     & Bloques try-catch \\
\bottomrule
\end{tabular}
\caption{Las 14 características extraídas del código Java}
\end{table}

Las características de conteo \texttt{code\_length}, \texttt{token\_count}, \texttt{brace\_count} y \texttt{parenthesis\_count} se transforman mediante $\log(1+x)$ para reducir el sesgo de distribución.

\subsection{Resultados del entrenamiento}

El modelo fue entrenado sobre 42,176 muestras (21,088 por clase) y evaluado sobre el 20\% reservado para prueba (8,436 muestras). Los resultados obtenidos superan ampliamente el umbral mínimo del 82\% exigido por la rúbrica.

\begin{table}[H]
\centering
\begin{tabular}{@{}lll@{}}
\toprule
\textbf{Métrica} & \textbf{Valor} & \textbf{Requisito} \\
\midrule
Accuracy (test set)         & \textbf{95.46\%} & $\geq$ 82\% \\
Precision (weighted)        & 0.9546           & ---          \\
Recall (weighted)           & 0.9546           & ---          \\
F1-Score (weighted)         & 0.9546           & ---          \\
Media CV 5-fold             & \textbf{95.34\%} $\pm$ 0.22\% & $\geq$ 82\% \\
\bottomrule
\end{tabular}
\caption{Métricas de evaluación del clasificador XGBoost}
\end{table}

\textbf{Reporte de clasificación por clase:}

\begin{table}[H]
\centering
\begin{tabular}{@{}lcccc@{}}
\toprule
\textbf{Clase} & \textbf{Precision} & \textbf{Recall} & \textbf{F1-score} & \textbf{Support} \\
\midrule
SEGURO     & 0.95 & 0.96 & 0.95 & 4,218 \\
VULNERABLE & 0.96 & 0.95 & 0.95 & 4,218 \\
\midrule
accuracy   & \multicolumn{3}{c}{0.95} & 8,436 \\
macro avg  & 0.95 & 0.95 & 0.95 & 8,436 \\
\bottomrule
\end{tabular}
\caption{Clasificación por clase en el conjunto de prueba}
\end{table}

\textbf{Scores de validación cruzada (5-fold):}
\[
[0.9517,\ 0.9576,\ 0.9514,\ 0.9527,\ 0.9536] \quad \bar{x} = 0.9534 \pm 0.0022
\]

\textbf{Matriz de confusión:}
\[
\begin{pmatrix}
4030 & 188 \\
195  & 4023
\end{pmatrix}
\quad
\begin{array}{l}
\text{TN (seguros correctos): 4,030} \\
\text{FP (seguros falseados): 188} \\
\text{FN (vulnerables no detectados): 195} \\
\text{TP (vulnerables detectados): 4,023}
\end{array}
\]

El número de falsos negativos (195) es comparable al de falsos positivos (188), lo que indica que el umbral de 0.35 está bien calibrado para el contexto de seguridad.

\subsection{Umbral de decisión}

Se aplica un umbral de probabilidad de \textbf{0.35} en lugar del estándar 0.50, priorizando la sensibilidad (recall) sobre la especificidad. Esto significa que si la probabilidad de que el código sea VULNERABLE supera el 35\%, el pipeline lo rechaza. Esta decisión se justifica en un contexto de seguridad donde los falsos negativos (código vulnerable que pasa) son más costosos que los falsos positivos (código seguro rechazado).

% -------------------------------------------------------
\section{Pipeline CI/CD}
% -------------------------------------------------------

\subsection{Job 1: Security Check}

\textbf{Trigger:} Pull Request hacia la rama \texttt{test}.

\begin{lstlisting}[style=yaml, caption={Extracto del job de análisis de seguridad}]
security-check:
  runs-on: ubuntu-latest
  if: github.event_name == 'pull_request' && github.base_ref == 'test'
  steps:
    - uses: actions/checkout@v3
      with:
        fetch-depth: 0
    - name: Telegram - Start Security Check
      run: curl -s -X POST "https://api.telegram.org/bot${{ secrets.TELEGRAM_TOKEN }}/sendMessage" ...
    - name: Analyze code
      run: python3 /tmp/analyze.py  # XGBoost predict_proba >= 0.35
    - name: Comentario en PR con resultado
      if: steps.analyze.outcome == 'failure'
      uses: actions/github-script@v6
    - name: Create Issue
      if: steps.analyze.outcome == 'failure'
      uses: actions/github-script@v6
\end{lstlisting}

Acciones al detectar código \textbf{VULNERABLE}:
\begin{itemize}
  \item El pipeline termina con código de salida 1 (bloqueo del merge).
  \item Se publica un comentario en el PR con el reporte detallado de features sospechosas.
  \item Se crea automáticamente una \textit{issue} en GitHub con las etiquetas \texttt{security-issue} y \texttt{fixing-required}.
  \item Se envía notificación inmediata por Telegram con el detalle de las vulnerabilidades.
\end{itemize}

\subsection{Job 2: Unit Tests}

\textbf{Dependencia:} \texttt{needs: security-check} — sólo se ejecuta si el código fue clasificado como SEGURO.

Se ejecutan en secuencia:
\begin{enumerate}
  \item \textbf{pytest} sobre el directorio \texttt{tests/} (25+ casos Python).
  \item \textbf{Maven} (\texttt{mvn test}) sobre los tests JUnit 5 de la aplicación Spring Boot (8 casos).
\end{enumerate}

En caso de fallo se crea una \textit{issue} con la etiqueta \texttt{tests-failed} y se envía notificación Telegram.

\subsection{Job 3: Merge y Despliegue}

\textbf{Trigger:} Push a la rama \texttt{test} (resultado del merge automático del PR aprobado).

\begin{enumerate}
  \item \texttt{git merge --no-ff origin/test} hacia \texttt{main}.
  \item \texttt{git push origin main}.
  \item \texttt{curl} al webhook de Render para activar el redeploy.
  \item Notificación Telegram con la URL de producción.
\end{enumerate}

% -------------------------------------------------------
\section{API Flask de Predicción}
% -------------------------------------------------------

La aplicación Flask (\texttt{app/app.py}) expone el modelo de ML como servicio REST, utilizado tanto por el pipeline como por el bot de Telegram y usuarios externos.

\begin{table}[H]
\centering
\begin{tabular}{@{}llp{7cm}@{}}
\toprule
\textbf{Endpoint}       & \textbf{Método} & \textbf{Descripción} \\
\midrule
\texttt{/health}        & GET             & Estado del servicio y disponibilidad del modelo \\
\texttt{/predict}       & POST            & Clasifica un fragmento de código Java \\
\texttt{/batch-predict} & POST            & Clasifica múltiples fragmentos en lote \\
\texttt{/model-info}    & GET             & Metadata del modelo (tipo, features, clases) \\
\bottomrule
\end{tabular}
\caption{Endpoints de la API Flask}
\end{table}

\textbf{URL de producción:} \url{https://secure-devops-pipeline-1.onrender.com}

El servidor de producción es Gunicorn, iniciado por el script \texttt{start.py} que también verifica la existencia del modelo y lo entrena si no está disponible.

% -------------------------------------------------------
\section{Aplicación Spring Boot}
% -------------------------------------------------------

La aplicación de demostración es una API REST CRUD de gestión de usuarios, desarrollada con Spring Boot 3.2.0 y Java 17. Sirve como código objeto del análisis de seguridad del pipeline.

\begin{itemize}
  \item \textbf{Entidad:} \texttt{Usuario} (id, nombre, email)
  \item \textbf{Validaciones:} \texttt{@NotBlank}, \texttt{@Email} mediante Bean Validation
  \item \textbf{Almacenamiento:} en memoria (\texttt{HashMap} + \texttt{AtomicLong})
  \item \textbf{Endpoints:} CRUD completo (\texttt{/usuarios}) y health check (\texttt{/health})
\end{itemize}

El frontend estático (\texttt{index.html}, \texttt{script.js}, \texttt{styles.css}) se sirve desde \texttt{src/main/resources/static/}.

% -------------------------------------------------------
\section{Notificaciones y Automatización}
% -------------------------------------------------------

\subsection{Bot de Telegram}

Se utiliza un bot propio de Telegram con token almacenado en GitHub Secrets (\texttt{TELEGRAM\_TOKEN}). El bot envía mensajes estructurados en los siguientes eventos:

\begin{enumerate}
  \item Inicio de análisis de seguridad
  \item Resultado SEGURO del clasificador
  \item Resultado VULNERABLE con detalle de features sospechosas
  \item Tests Python y Java aprobados
  \item Tests fallidos (con bloqueo)
  \item Merge \texttt{test}$\to$\texttt{main} exitoso
  \item Deploy exitoso en Render con URL
  \item Fallo en merge o deploy
\end{enumerate}

Adicionalmente, el archivo \texttt{telegram\_bot.py} implementa un bot interactivo que permite a cualquier usuario enviar código Java y recibir la predicción del modelo directamente en la conversación de Telegram.

\subsection{Issues automáticas en GitHub}

El workflow crea \textit{issues} en el repositorio de forma automática ante dos condiciones:

\begin{itemize}
  \item \textbf{Código vulnerable:} etiquetas \texttt{security-issue} + \texttt{fixing-required}, con el reporte completo del análisis.
  \item \textbf{Tests fallidos:} etiqueta \texttt{tests-failed}, con referencia al PR y al log del pipeline.
\end{itemize}

% -------------------------------------------------------
\section{Despliegue en Producción}
% -------------------------------------------------------

\subsection{Plataforma: Render}

Se utiliza Render como proveedor de PaaS gratuito. La configuración de despliegue está definida en el archivo \texttt{Procfile}:

\begin{lstlisting}[style=yaml]
web: python start.py
\end{lstlisting}

El script \texttt{start.py} realiza las siguientes acciones en el arranque:
\begin{enumerate}
  \item Crea los directorios necesarios (\texttt{model/}, \texttt{data/raw/}, \texttt{data/processed/}).
  \item Verifica si existe \texttt{model/classifier\_model.joblib}; si no existe, ejecuta el entrenamiento.
  \item Inicia el servidor Gunicorn con la aplicación Flask.
\end{enumerate}

El deploy se activa automáticamente mediante un webhook de Render llamado desde el \textit{job} de CI/CD.

\subsection{Docker}

Se incluye un \texttt{Dockerfile} multi-stage:
\begin{itemize}
  \item \textbf{Etapa de build:} \texttt{eclipse-temurin:17-jdk-alpine} con Maven para compilar el proyecto Java.
  \item \textbf{Etapa de runtime:} \texttt{eclipse-temurin:17-jre-alpine} (imagen mínima) que expone el puerto 8080.
\end{itemize}

% -------------------------------------------------------
\section{Pruebas}
% -------------------------------------------------------

\subsection{Pruebas Python (pytest)}

\begin{table}[H]
\centering
\begin{tabular}{@{}ll@{}}
\toprule
\textbf{Archivo}                   & \textbf{Casos de prueba} \\
\midrule
\texttt{tests/test\_app.py}        & Importación Flask, endpoint /health, carga del modelo \\
\texttt{tests/test\_feature\_extractor.py} & 18 casos: código vacío, funciones peligrosas, SQL, \\
                                    & sanitización, excepciones, depth, array, 14 features \\
\bottomrule
\end{tabular}
\caption{Suite de pruebas Python}
\end{table}

\subsection{Pruebas Java (JUnit 5)}

La clase \texttt{UsuarioServiceTest.java} contiene 8 casos de prueba que verifican las operaciones CRUD completas del servicio de usuarios, incluyendo casos límite como IDs inexistentes.

% -------------------------------------------------------
\section{Estructura del Repositorio}
% -------------------------------------------------------

\begin{verbatim}
secure-devops-pipeline/
├── .github/workflows/main.yml     -- pipeline completo (3 jobs)
├── model/
│   ├── train_model.py             -- script de entrenamiento XGBoost
│   ├── classifier_model.joblib    -- modelo serializado
│   ├── scaler.joblib              -- StandardScaler serializado
│   └── feature_names.json         -- orden de las 14 features
├── security_checker/
│   └── feature_extractor.py      -- JavaCodeAnalyzer (14 features)
├── app/
│   └── app.py                    -- API Flask (4 endpoints)
├── notebooks/
│   ├── 01_EDA_Vulnerability_Dataset.ipynb
│   ├── 02_Data_Preprocessing.ipynb
│   └── 03_Create_Classification_Dataset.ipynb
├── src/                           -- Spring Boot 3.2.0 (Java 17)
│   ├── main/java/com/ejemplo/usuarios/
│   └── test/java/com/ejemplo/usuarios/
├── tests/                         -- pytest (25+ casos)
├── Dockerfile                     -- multi-stage build
├── Procfile                       -- web: python start.py
├── start.py                       -- entrypoint Render
├── telegram_bot.py                -- bot interactivo Telegram
├── requirements.txt               -- 14 dependencias Python
└── pom.xml                        -- Spring Boot 3.2.0, Java 17
\end{verbatim}

% -------------------------------------------------------
\section{Criterios de Evaluación y Cumplimiento}
% -------------------------------------------------------

\begin{table}[H]
\centering
\begin{tabular}{@{}p{9cm}cc@{}}
\toprule
\textbf{Criterio}                                              & \textbf{Puntaje} & \textbf{Estado} \\
\midrule
Funcionalidad completa del pipeline (automatización total)     & 6 pts  & \checkmark \\
Modelo de minería de datos propio y efectivo (no LLM)          & 6 pts  & \checkmark \\
Notificaciones Telegram en todas las fases + issues automáticas & 3 pts  & \checkmark \\
Despliegue automático en proveedor gratuito y funcional        & 3 pts  & \checkmark \\
Calidad del informe y documentación (README + notebook)        & 2 pts  & \checkmark \\
\midrule
\textbf{Total}                                                 & \textbf{20 pts} & \\
\bottomrule
\end{tabular}
\caption{Cumplimiento de criterios de evaluación}
\end{table}

\begin{table}[H]
\centering
\begin{tabular}{@{}p{9cm}c@{}}
\toprule
\textbf{Requisito específico}                                  & \textbf{Cumplido} \\
\midrule
Modelo ML clásico (no LLM), archivo \texttt{.joblib} entregado & \checkmark \\
Dataset público con $>$82\% accuracy en validación cruzada 5-fold & \checkmark \\
Features mínimas: tokens, AST depth, funciones peligrosas, sanitización & \checkmark \\
Bot Telegram propio con token en GitHub Secrets               & \checkmark \\
Despliegue real y funcional (Render, online)                  & \checkmark \\
Branch protection en \texttt{test} y \texttt{main}           & \checkmark \\
PR bloqueado automáticamente si VULNERABLE                    & \checkmark \\
Comentario detallado en PR + issue automática                 & \checkmark \\
Merge automático \texttt{dev}$\to$\texttt{test}$\to$\texttt{main} & \checkmark \\
Notebooks de EDA e incluidos en el repositorio                & \checkmark \\
\bottomrule
\end{tabular}
\caption{Checklist de requisitos del proyecto}
\end{table}

% -------------------------------------------------------
\section{Conclusiones}
% -------------------------------------------------------

Se diseñó e implementó exitosamente un pipeline CI/CD seguro que cumple todos los requisitos del proyecto integrador. Las conclusiones principales son:

\begin{enumerate}
  \item \textbf{Automatización total:} el flujo \texttt{dev}$\to$\texttt{test}$\to$\texttt{main} y el despliegue en Render se ejecutan sin intervención humana, garantizando consistencia y trazabilidad.
  \item \textbf{ML clásico efectivo:} XGBoost con 14 features léxicas y estructurales del código Java ofrece una solución interpretable, liviana y desplegable en entornos de CI/CD sin requerir GPU ni infraestructura especializada.
  \item \textbf{Umbral conservador:} el umbral de 0.35 para clasificación VULNERABLE prioriza la seguridad sobre la conveniencia del desarrollador, alineándose con los principios de Shift-Left Security.
  \item \textbf{Integración completa:} la combinación de GitHub Actions, Telegram, Render y GitHub Issues crea un ciclo de retroalimentación inmediato y completamente automatizado.
  \item \textbf{Extensibilidad:} la arquitectura modular (extractor de features separado del clasificador) permite añadir nuevas features o cambiar el algoritmo ML sin modificar el pipeline.
\end{enumerate}

% -------------------------------------------------------
\section*{Referencias}
% -------------------------------------------------------

\begin{enumerate}[label={[\arabic*]}]
  \item Chen, T., \& Guestrin, C. (2016). XGBoost: A scalable tree boosting system. \textit{Proceedings of the 22nd ACM SIGKDD International Conference on Knowledge Discovery and Data Mining}, 785--794.
  \item Scikit-learn developers. (2023). \textit{scikit-learn: Machine Learning in Python}. \url{https://scikit-learn.org}
  \item DevOps Research and Assessment (DORA). (2023). \textit{Accelerate State of DevOps Report}. Google Cloud.
  \item GitHub. (2024). \textit{GitHub Actions Documentation}. \url{https://docs.github.com/en/actions}
  \item Render. (2024). \textit{Render Platform Documentation}. \url{https://render.com/docs}
  \item OWASP Foundation. (2021). \textit{OWASP Top Ten}. \url{https://owasp.org/www-project-top-ten/}
  \item Kim, G., Humble, J., Debois, P., \& Willis, J. (2016). \textit{The DevOps Handbook}. IT Revolution Press.
\end{enumerate}

\end{document}
